\chapter{Experimental Setup}\label{ch:setup}
\section{Tasks and Environments}\label{sec:setup:tasks}
\section{Datasets}\label{sec:setup:data}
\section{Baselines}\label{sec:setup:baselines}
\section{Metrics}\label{sec:setup:metrics}
\section{Training and Evaluation Protocol}\label{sec:setup:protocol}

% -----------------------------------------------------------------------------
% EXCEPTION to the "bone only" rule of Chapters 5-8 above. The compute
% environment is NOT blocked on the run queue: it is a property of the machine,
% settled and independently verifiable today (sinfo/scontrol/lscpu, 2026-09-09,
% plus every sbatch log header). Written now so it does not have to be
% reconstructed later. Source of record, with the raw command output:
%   Writing/Auxiliary/Methodology_Sources/AUX_compute_environment.md
% Everything else in Chapters 5-8 stays bone.
% -----------------------------------------------------------------------------

\subsection*{Compute environment}

All training and evaluation reported in this thesis ran on a single shared GPU node of the
institute's cluster, scheduled with SLURM. The node carries two 16-core AMD EPYC~7282 processors
(32~physical cores, 64~hardware threads) and eight NVIDIA RTX~A5000 accelerators with 24~GB of
device memory each. \textbf{Every job was allocated exactly one GPU}: no result in this thesis comes
from a distributed or multi-GPU run. The full hardware and software inventory, including package
versions, is given in \autoref{app:repro}, \autoref{tab:compute}.
\srcnote{Slurm\_Codes/sbatch/templates/2026\_04\_30\_job\_template.sh; every log header in Slurm\_Codes/logs/}

Two properties of this machine bear directly on how the wall-clock measurements in
\autoref{ch:results} may be read, and both argue for reading them as \emph{indicative} rather than as
a hardware benchmark.

\emph{The node is shared.} It serves the institute's student and staff GPU queues simultaneously,
and a snapshot taken while preparing this section found six of the eight accelerators and 42 of the
64 threads already committed to other users' jobs. What this does \emph{not} affect is the
accelerator itself: the scheduler neither oversubscribes GPUs nor preempts running jobs, so a job
holds its device exclusively for its lifetime. Contention is therefore confined to the CPU, the
memory system and I/O.

\emph{The CPU-to-GPU ratio is low, and this work is CPU-bound.} Thirty-two physical cores serve
eight accelerators, so a fully occupied node offers roughly four physical cores per GPU. Both halves
of the evaluation loop --- the MuJoCo rollout and the per-step projection, which is solved by SLSQP
on the CPU (\autoref{sec:method:proj}) --- run on those cores rather than on the accelerator.
Evaluation throughput on this machine is thus plausibly limited by the CPU side and not by the
RTX~A5000. This holds independently of contention, that is, it would hold on an idle node as well;
it is the more fundamental of the two caveats.

Both observations are recorded here rather than in \autoref{sec:disc:threats} because they concern
the measurement apparatus rather than the validity of a claim. Their consequence for the timing
figures specifically is taken up in \autoref{sec:res:summary}.

